% !TeX spellcheck = it_IT

\documentclass[handout]{beamer}

\usepackage[utf8]{inputenc}
\usepackage[italian]{babel}
\usepackage{multicol}
\usetheme{metropolis}
\usepackage{forest}

\usepackage{listings}
\usepackage{color} %red, green, blue, yellow, cyan, magenta, black, white
\definecolor{mygreen}{RGB}{28,172,0} % color values Red, Green, Blue
\definecolor{mylilas}{RGB}{170,55,241}
\lstset{language=Matlab,%
	%basicstyle=\color{red},
	breaklines=true,%
	morekeywords={matlab2tikz},
	keywordstyle=\color{blue},%
	morekeywords=[2]{1}, keywordstyle=[2]{\color{black}},
	identifierstyle=\color{black},%
	stringstyle=\color{mylilas},
	commentstyle=\color{mygreen},%
	showstringspaces=false,%without this there will be a symbol in the places where there is a space
	numbers=left,%
	numberstyle={\tiny \color{black}},% size of the numbers
	numbersep=9pt, % this defines how far the numbers are from the text
	emph=[1]{for,end,break},emphstyle=[1]\color{red}, %some words to emphasise
	%emph=[2]{word1,word2}, emphstyle=[2]{style},
	tabsize=2,
}



\usepackage{tikz}
\usetikzlibrary{shapes,arrows,positioning}

\tikzstyle{startstop} = [ellipse, draw, fill=bg, text=structure, text width=3em, text badly centered, inner sep=0pt]
\tikzstyle{block} = [rectangle, draw, fill=bg, text=structure, text width=5em, text centered]
\tikzstyle{decision} = [diamond, aspect=2, draw, fill=bg, text=structure, text width=3em, text badly centered, inner sep=0pt]
\tikzstyle{input} = [rectangle, draw, fill=bg, text=structure, text width=4em, text centered, rounded rectangle, rounded rectangle east arc=none]
\tikzstyle{output} = [rectangle, draw, fill=bg, text=structure, text width=4em, text centered, rounded rectangle, rounded rectangle west arc=none]

\tikzstyle{line} = [draw, -latex']


\title{Informatica B - Esercitazione 10}
\subtitle{Matlab - Diagrammi e ripasso}

\author{Stefano Cereda\\
		stefano1.cereda@mail.polimi.it
	}
\date{12/12/2017}
\institute[PoliMi]{\vspace{0.5cm}\centering Politecnico di Milano \\ \vspace{0.2cm}
	\includegraphics[width=0.2\linewidth]{./logopolimi}}
%\logo{\includegraphics[width=15mm]{./logopolimi}}

\setbeamercovered{invisible}

\makeindex

\begin{document}
\begin{frame}
	\maketitle
\end{frame}

\begin{frame}[allowframebreaks]{Esercizio 1 - Caldaia}
	In un cinema di Milano la temperatura nella sala è regolata in modo
	automatico; un sensore monitora la temperatura rilevando un valore ogni
	minuto. I valori letti vengono inseriti in un array temp.
	
	La caldaia si accende solo quando la temperatura rilevata é inferiore ad una
	certa soglia e si spegne non appena viene raggiunta tale soglia.
	\begin{enumerate}
		\item  Scrivere in Matlab uno script che:
		\begin{enumerate}
			\item Generi in maniera random un vettore di temperature e che chieda
			in ingresso un valore soglia ed un parametro ``costoAlMinuto''
			indicante il costo al minuto del gas consumato dalla caldaia;
			\item Calcoli e stampi il costo totale del gas consumato ed un vettore
			contenenti gli indici dell’array temp corrispondenti ai minuti nei
			quali la caldaia é rimasta accesa.
			
			Per il calcolo del costo totale del gas consumato si consideri che se il
			totale dei minuti nei quali la caldaia ha funzionato supera i 30 minuti,
			allora il costo al minuto va diminuito del 20\%.
		\end{enumerate}
		\item Disegnare il grafico della temperatura al variare del tempo, evidenziando
		sul grafico stesso:
		\begin{enumerate}
			\item Il valore soglia con una retta orizzontale nera a doppio spessore
			\item I minuti nei quali la caldaia ha funzionato, marcandoli con il simbolo
			‘X’ rosso (senza congiungerli).
		\end{enumerate}
	\end{enumerate}
\end{frame} 

\begin{frame}[allowframebreaks, fragile]{Esercizio 2 - Roulette}
	Nel gioco della roulette si ha una ruota raffigurante 37
	numeri da 0 a 36. Ad ogni turno un giocatore può
	scommettere su un numero o su particolari combinazioni
	degli stessi. Alla fine delle puntate la ruota selezionerà il
	numero vincente.
	
	\begin{enumerate}
		\item Si scriva una funzione che simuli la ruota e generi il
		numero vincente utilizzando la funzione \emph{rand}.
		\item Si scriva una funzione che provi ad eseguire N giri e salvi
		i risultati delle frequenze di ogni numero in un array di
		37 elementi (il primo elemento per lo zero, il secondo per
		l'uno e così via).
		\item Si tracci il grafico delle frequenze prima utilizzando
		\emph{rand} come generatore e poi utilizzando la funzione
		\emph{rand\_truccata} mostrata di seguito e commentare il
		risultato.
		\item Qual è la probabilità di vincere se si scommette
		sull'uscita di un numero dispari? Scrivere una funzione
		che calcoli questa probabilità effettuando 100000
		simulazioni di giocate.
	\end{enumerate}
	
	\framebreak
	\begin{lstlisting}
function [n]=rand_truccata()
	r=rand();
	if r<0.02	
		n=0;
	else
		n=rand();
	end
	\end{lstlisting}
\end{frame}

\begin{frame}[allowframebreaks]{Esercizio 3 - Dado - Lasciare da fare}
	Nel gioco dei dadi ad ogni faccia del dado è associato un
	numero da 1 a 6.
	
	\begin{enumerate}
		\item Si scriva una funzione che simuli il dado e generi il
		numero vincente utilizzando la funzione rand.
		\item Si scriva una funzione che provi ad eseguire N lanci e
		salvi i risultati delle frequenze di ogni numero in un array
		di 6 elementi (il primo elemento per l’uno, il secondo per
		il due e così via).
		\item Qual è la probabilità di vincere se si scommette
		sull'uscita di un numero? Scrivere una funzione che
		calcoli questa probabilità effettuando 100000 simulazioni
		di lanci.
	\end{enumerate}
\end{frame}

\begin{frame}[allowframebreaks]{Esercizio 4 - Prova in itinere 2014}
	Un vettore x si definisce vettore triplo di un altro vettore y quando x e
	y sono di uguale lunghezza e ogni elemento di x è maggiore del triplo
	del corrispondente elemento di y.
	
	Scrivere in Matlab/Octave una funzione \alert{ricorsiva} \emph{vettoreTriplo} che
	riceve in ingresso due vettori \alert{colonna} \emph{a} e \emph{b}. La funzione
	vettoreTriplo:
	\begin{enumerate}
		\item Restituisce -1 se \emph{a} e \emph{b} hanno lunghezza diversa, o se non sono due vettori colonna.
		\item Restituisce 0 se \emph{a} e \emph{b} sono due vettori colonna di uguale
		lunghezza, ma non è vero che ogni elemento di \emph{a} è maggiore
		del triplo del corrispondente elemento di \emph{b}.
		\item Restituisce 1 se \emph{a} e \emph{b} sono due vettori colonna di uguale
		lunghezza e ogni elemento di \emph{a} è maggiore
		del triplo del corrispondente elemento di \emph{b}.
	\end{enumerate}

	Esempi:
	\begin{itemize}
		\item vettoreTriplo ([1; 2; 3; 4], [2 5 1 10]) restituisce -1
		\item vettoreTriplo ( [1 2 3 4; 5 3 5 5], [2; 2; 3; 9]) restituisce -1
		\item vettoreTriplo ([7; 2; 4; 4], [2; 5; 1; 10]) restituisce 0
		\item vettoreTriplo([7; 16; 40; 31], [2; 5; 13; 10]) restituisce 1
	\end{itemize}

	Nell’implementazione \alert{non} è consentito utilizzare cicli, e quindi \alert{non} si possono usare le istruzioni \emph{while} e \emph{for}. Inoltre, \alert{non} è possibile utilizzare alcuna funzione di libreria disponibile in Matlab/Octave ad eccezione di \emph{length} e \emph{size}.
\end{frame}

\begin{frame}[fragile]{Esercizio 5 - Paraboloide}
	Si disegni la superficie del paraboloide iperbolico dato dalla seguente
	equazione:
	\begin{lstlisting}
z=x.^2-y.^2
	\end{lstlisting}

	per valori di x e y nell'intervallo [-5,+5].\\
	Impostare titoli ed etichette.
\end{frame}

\begin{frame}[fragile, allowframebreaks]{Esercizio 6 - Funzione misteriosa TDE 17/09/2012}
	Si consideri la seguente funzione ricorsiva in linguaggio MATLAB:
	\begin{lstlisting}
function n = calcola(x)
	if (floor(x/10) == 0)
		n = 1;
	else
		n = 1 + calcola(floor(x/10));
end
\end{lstlisting}
	
	Rispondere ai seguenti quesiti.
	\begin{enumerate}
		\item Simulare l'esecuzione della funzione e indicare il valore restituito in corrispondenza di ciascuna delle
		chiamate: calcola(4) e calcola(123)
		\item 
		Assumendo che la funzione calcola riceva come parametro un qualsiasi numero intero positivo n , dire qual
		è il risultato dell'esecuzione di calcola(n) al variare di n. Motivare adeguatamente la risposta.
	\end{enumerate}

\alert{Nota:} assieme all'esercizio sulle fantasquadre (esercitazione 5, 10 punti) e all'esercizio sugli iban (esercitazione 8, 10 punti) questo esercizio da 6 punti completa il tema d'esame del 17 Settembre 2012.
\end{frame}

\end{document}